

\newcommand{\ltree}[4]{%
\begin{tikzpicture}[
    every node/.style={font=\small},
    circle node/.style={circle, draw, inner sep=1pt,thick},
    triangle node/.style={
        regular polygon,
        regular polygon sides=3,
        draw=black,
        thick,
        inner sep=-0.5pt, 
%        minimum size=2mm
    },
    level distance=0.7cm,
    sibling distance=0.7cm,
    % Force all edges to go to the top of the child node
        edge from parent/.style={ultra thick},
    edge from parent path={
        (\tikzparentnode.south) -- (\tikzchildnode.north)
    }
]
%
\node[circle node] (x) {$x$}
    child {
        node[triangle node] (A) {$A$}
        edge from parent[draw=#1]
    }
    child {
        node[circle node] (y) {$y$}
        edge from parent[draw=#2]
        child {
            node[triangle node] (B) {$B$}
            edge from parent[draw=#3]
        }
        child {
            node[triangle node] (C) {$C$}
            edge from parent[draw=#4]
        }
    };
%
\end{tikzpicture}
}


\newcommand{\rtree}[4]{%
\begin{tikzpicture}[
    every node/.style={font=\small},
    circle node/.style={circle, draw, inner sep=1pt,thick},
    triangle node/.style={
        regular polygon,
        regular polygon sides=3,
        draw=black,
        inner sep=-0.5pt, 
        thick
%        minimum size=2mm
    },
    level distance=0.7cm,
    sibling distance=0.7cm,
    % Force all edges to go to the top of the child node
    edge from parent/.style={ultra thick},
    edge from parent path={
        (\tikzparentnode.south) -- (\tikzchildnode.north)
    }
]
%
\node[circle node] (x) {$y$}
    child {
        node[circle node] (y) {$x$}
        edge from parent[draw=#1]
        child {
            node[triangle node] (A) {$A$}
            edge from parent[draw=#3]
        }
        child {
            node[triangle node] (B) {$B$}
            edge from parent[draw=#4]
        }
        }
    child {
        node[triangle node] (C) {$C$}
        edge from parent[draw=#2]
    };
%
\end{tikzpicture}
}



\newcommand{\lta}{\ltree{red}{black}{red}{black}}
\newcommand{\ltb}{\ltree{red}{black}{black}{blue}}
\newcommand{\ltc}{\ltree{black}{blue}{red}{black}}
\newcommand{\ltd}{\ltree{black}{blue}{black}{blue}}

\newcommand{\rta}{\rtree{red}{black}{red}{black}}
\newcommand{\rtb}{\rtree{red}{black}{black}{blue}}
\newcommand{\rtc}{\rtree{black}{blue}{red}{black}}
\newcommand{\rtd}{\rtree{black}{blue}{black}{blue}}

\newcounter{casenum}

\newcommand{\ra}{%
{\stepcounter{casenum}%
\raisebox{42pt}{$\mathclap{\text{Case \Roman{casenum}}}$}}%
{\raisebox{20pt}{$\mathclap{\Rightarrow}$}}%
}

\newcommand{\lhc}[1]{Left heavy change: #1}



\begin{figure}
\resizebox{\textwidth}{!}{%
\begin{tabular}{|c|c|c|c|}
\hline
\lta \textcolor{\cbc}{\ra} \rta &
\lta \textcolor{\cac}{\ra}  \rtb &
\lta \textcolor{\cac}{\ra}  \rtc &
\lta \textcolor{\cac}{\ra}  \rtd 
\\
\lhc{0}
&
\lhc{-1}
&
\lhc{-1}
&
\lhc{-2}
\\ \hline
\ltb \textcolor{\cac}{\ra}  \rta &
\ltb \textcolor{\ccc}{\ra}  \rtb &
\ltb \textcolor{\ccc}{\ra}  \rtc &
\ltb \textcolor{\cac}{\ra}  \rtd 
\\
\lhc{+1}
&
Impossible
&
Impossible
&
\lhc{-1}
\\ \hline
\ltc \textcolor{\cac}{\ra}  \rta &
\ltc \textcolor{\cdc}{\ra}  \rtb &
\ltc \textcolor{\ccc}{\ra}  \rtc &
\ltc \textcolor{\cac}{\ra}  \rtd 
\\
\lhc{+1}
&
\lhc{0}
&
Impossible
&
\lhc{-1}
\\ \hline
\ltd \textcolor{\cac}{\ra} \rta &
\ltd \textcolor{\cac}{\ra}  \rtb &
\ltd \textcolor{\cac}{\ra}  \rtc &
\ltd \textcolor{\cbc}{\ra}  \rtd 
\\
\lhc{+2}
&
\lhc{+1}
&
\lhc{+1}
&
\lhc{0}
\\ \hline 
\end{tabular}
}
\caption{
The cases. We consider a left rotation on $x$ with right child $y$. A red edge indicates a left heavy node above and a blue edge indicates a right heavy node above. We consider all 16 possible cases of left and right heaviness for $x$ and $y$ before and after the rotation. 
In \textcolor{\cac}{Cases II-V, VIII, IX, XII-XV} the number of left heavy nodes change by 1 or 2 and thus the left color changes. 
In \textcolor{\cdc}{Case X} the only change in heavy paths is the relative ordering of $x$ and $y$ is swapped, each retaining their heaviness; thus the path color changes.
\textcolor{\ccc}{Cases VI, VII and XI} are impossible. 
\textcolor{\cbc}{Cases I and XVI} are possible but have no change in the number of left heavy nodes and thus require a more detailed argument and are the source of the logarithmic number of colors.
}
\label{f1}
\end{figure}
